\begin{definition} \label{definition:real_power_series_and_radius_of_convergence}
    \TODO{AI generated, verify}
    Let $(a_n)_{n=0}^\infty$ be a sequence of real numbers, and let $c \in \mathbb{R}$. The expression
    $$\sum_{n=0}^\infty a_n (x - c)^n$$
    is called a \hldef{real power series centered at $c$}, where the coefficients are denoted by \hl{$a_n$}. For each $x \in \mathbb{R}$, the \hldef{$N$th partial sum of the power series at $x$}, denoted by \hl{$S_N(x)$}, is defined as
    $$S_N(x) := \sum_{n=0}^N a_n (x - c)^n.$$

    \TODO{cauchy hadamard formula}
    The \hldef{radius of convergence} of the power series, denoted by \hl{$R$}, is the nonnegative extended real number given by the Cauchy--Hadamard formula:
    $$R := \frac{1}{\limsup_{n \to \infty} |a_n|^{1/n}}$$
    with the conventions $1/0 = \infty$ and $1/\infty = 0$. The power series \CrefAndHyperrefIfExist{definition:absolute_convergence_of_a_series_in_module_over_a_commutative_ring_with_an_absolute_value_valued_in_an_ordered_semiring}{converges absolutely} for all $x \in \mathbb{R}$ with $|x - c| < R$, and diverges for all $x$ with $|x - c| > R$. The open interval $(c - R, c + R)$ is called the \hldef{interval of convergence}.
\end{definition}
